% ============================================================================
%  Rozdział: Interpretacja parametrów intrinsics dla Intel RealSense
% ============================================================================
\documentclass[12pt,a4paper]{report}
\usepackage[utf8]{inputenc}
\usepackage[T1]{fontenc}
\usepackage[polish]{babel}
\usepackage{amsmath,amssymb}
\usepackage{booktabs}
\usepackage{array}
\usepackage{hyperref}
\usepackage{geometry}
\geometry{margin=2.5cm}

\begin{document}

\chapter{Interpretacja parametrów \texttt{intrinsics} zapisanych z kamery Intel RealSense}
\label{chap:intrinsics-realsense}

\section{Wprowadzenie}
Funkcja \texttt{intrinsics\_to\_dict(intr)} zapisuje do pliku JSON podstawowe parametry wewnętrzne kamery:
\texttt{width}, \texttt{height}, \texttt{fx}, \texttt{fy}, \texttt{cx}, \texttt{cy} oraz strukturę
\texttt{distortion}. Parametry te pochodzą bezpośrednio z obiektu \texttt{rs2\_intrinsics}
udostępnianego przez bibliotekę Intel RealSense SDK 2.0. Dokumentacja API producenta
precyzuje znaczenie każdego pola, w tym jego interpretację geometryczną i jednostki.

W notacji stosowanej w literaturze kalibracji kamery oraz w OpenCV parametry te tworzą
macierz wewnętrzną kamery
\begin{equation}
K =
\begin{bmatrix}
f_x & 0   & c_x \\
0   & f_y & c_y \\
0   & 0   & 1
\end{bmatrix}.
\end{equation}
Dokumentacja OpenCV wprost podaje, że $f_x$ i $f_y$ są wyrażone w jednostkach pikselowych,
a $(c_x, c_y)$ jest punktem głównym obrazu.

\section{Znaczenie poszczególnych pól}
Tabela~\ref{tab:intrinsics-meaning} zestawia pola zapisane przez funkcję oraz ich znaczenie
zgodne z dokumentacją \texttt{rs2\_intrinsics}.

\begin{table}[htbp]
    \centering
    \caption{Znaczenie pól zapisywanych przez \texttt{intrinsics\_to\_dict}.}
    \label{tab:intrinsics-meaning}
    \begin{tabular}{@{}p{2.6cm}p{3.2cm}p{8cm}@{}}
        \toprule
        \textbf{Pole} & \textbf{Jednostka} & \textbf{Znaczenie} \\
        \midrule
        \texttt{width} & piksele & Szerokość obrazu w pikselach. Dokumentacja \texttt{rs2\_intrinsics} definiuje to pole jako \emph{Width of the image in pixels}. \\
        \texttt{height} & piksele & Wysokość obrazu w pikselach. Dokumentacja definiuje to pole jako \emph{Height of the image in pixels}. \\
        \texttt{fx} & piksele & Ogniskowa w osi poziomej. Producent opisuje ją jako \emph{Focal length of the image plane, as a multiple of pixel width}, co odpowiada ogniskowej wyrażonej w skali piksela wzdłuż osi $x$. \\
        \texttt{fy} & piksele & Ogniskowa w osi pionowej. Producent opisuje ją jako \emph{Focal length of the image plane, as a multiple of pixel height}, czyli ogniskową w skali piksela wzdłuż osi $y$. \\
        \texttt{cx} & piksele & Współrzędna pozioma punktu głównego obrazu. W kodzie jest to \texttt{intr.ppx}; dokumentacja określa to pole jako \emph{Horizontal coordinate of the principal point of the image, as a pixel offset from the left edge}. \\
        \texttt{cy} & piksele & Współrzędna pionowa punktu głównego obrazu. W kodzie jest to \texttt{intr.ppy}; dokumentacja określa to pole jako \emph{Vertical coordinate of the principal point of the image, as a pixel offset from the top edge}. \\
        \texttt{distortion.model} & typ modelu & Nazwa modelu dystorsji przypisanego do danego strumienia obrazu. \\
        \texttt{distortion.coefficients} & bezwymiarowe & Pięcioelementowy wektor współczynników dystorsji. Dla modelu Brown--Conrady kolejność według dokumentacji RealSense to \mbox{$[k_1, k_2, p_1, p_2, k_3]$}. \\
        \bottomrule
    \end{tabular}
\end{table}

\section{Interpretacja geometryczna}
Parametry $f_x$ i $f_y$ nie są tutaj wyrażone w milimetrach, lecz w pikselach. Jest to bardzo
istotne przy interpretacji wyników, ponieważ w zapisie SDK RealSense nie dostajemy bezpośrednio
ogniskowej fizycznej soczewki, tylko ogniskową przeskalowaną do siatki obrazu. Z tego powodu
wartości \texttt{fx} i \texttt{fy} określają skalę projekcji perspektywicznej bezpośrednio
w układzie obrazu. Informacja, że są to jednostki pikselowe, nie jest założeniem interpretacyjnym,
lecz wynika wprost z dokumentacji producenta oraz z dokumentacji OpenCV.

Punkt główny $(c_x, c_y)$ określa położenie przecięcia osi optycznej z płaszczyzną obrazu.
W dokumentacji producenta jest on podany jako przesunięcie pikselowe od lewego oraz górnego
brzegu obrazu. Oznacza to, że wartości \texttt{cx} i \texttt{cy} należy interpretować w tym
samym układzie współrzędnych obrazu, w którym indeksowane są piksele. Dla obrazu o szerokości
$W$ i wysokości $H$ punkt główny idealnie centralny miałby współrzędne zbliżone do
$(W/2, H/2)$, lecz w praktyce może być od środka nieznacznie przesunięty.

\section{Znaczenie pola \texttt{distortion}}
Pole \texttt{distortion} zapisuje dwa elementy: nazwę modelu oraz wektor współczynników.
Dokumentacja RealSense dla struktury \texttt{rs2\_intrinsics} podaje, że dla modelu
Brown--Conrady współczynniki są uporządkowane jako
\begin{equation}
    [k_1, k_2, p_1, p_2, k_3],
\end{equation}
co jest zgodne z konwencją stosowaną również w OpenCV. W tej interpretacji:
\begin{itemize}
    \item $k_1$, $k_2$, $k_3$ opisują dystorsję radialną,
    \item $p_1$, $p_2$ opisują dystorsję tangencjalną.
\end{itemize}

Jest to ważne, ponieważ kolejność współczynników nie jest dowolna. Jeżeli dalsza analiza,
kalibracja lub porównanie z OpenCV zakłada model Brown--Conrady, to wektor z pliku JSON można
bezpośrednio odczytywać jako $[k_1, k_2, p_1, p_2, k_3]$.

\section{Do jakich dokumentów odnosi się ta interpretacja}
Znaczenie tych parametrów nie wynika z przypuszczeń ani z nieformalnych opisów w internecie,
tylko z dokumentacji bibliotek, które definiują używane typy danych i model kamery.
Najważniejsze odniesienia są następujące:
\begin{itemize}
    \item dokumentacja struktury \texttt{rs2\_intrinsics} w \texttt{librealsense2}, ponieważ to ona
        definiuje pola \texttt{width}, \texttt{height}, \texttt{ppx}, \texttt{ppy}, \texttt{fx}, \texttt{fy},
        \texttt{model} i \texttt{coeffs},
    \item dokumentacja enumeracji \texttt{rs2\_distortion}, ponieważ to ona opisuje, jak należy rozumieć
        nazwy modeli dystorsji przypisane do pola \texttt{model},
    \item dokumentacja OpenCV \texttt{calib3d}, ponieważ potwierdza standardową postać macierzy kamery
        oraz użycie parametrów $f_x$, $f_y$, $c_x$, $c_y$ i współczynników dystorsji
        $[k_1, k_2, p_1, p_2, k_3]$ w powszechnie stosowanej praktyce kalibracyjnej.
\end{itemize}

Z punktu widzenia poprawności biblioteki najistotniejsze jest to, że pola odczytywane w Pythonie
przez \texttt{pyrealsense2} odpowiadają bezpośrednio polom struktury \texttt{rs2\_intrinsics}
opisanej w dokumentacji \texttt{librealsense2}. Oznacza to, że skrypt nie tworzy własnej,
nieudokumentowanej interpretacji parametrów, lecz tylko serializuje dane zwracane przez oficjalne
API producenta.

\section{Dlaczego można uznać te informacje za wiarygodne}
Wiarygodność tej interpretacji opiera się na zgodności trzech poziomów opisu:
\begin{itemize}
    \item poziomu implementacyjnego: skrypt odczytuje pola bezpośrednio z obiektu zwracanego przez SDK,
    \item poziomu API: dokumentacja \texttt{rs2\_intrinsics} definiuje te pola jednoznacznie,
    \item poziomu modelu matematycznego: OpenCV używa tej samej notacji i tej samej interpretacji
        macierzy wewnętrznej kamery oraz wektora dystorsji.
\end{itemize}

W praktyce oznacza to, że gdy dokumentacja producenta mówi o \emph{pixel offset from the left edge},
to \texttt{cx} i \texttt{cy} należy interpretować jako współrzędne punktu głównego w układzie
obrazu, a gdy mówi o \emph{focal length as a multiple of pixel width/height}, to \texttt{fx}
i \texttt{fy} należy traktować jako ogniskowe wyrażone w pikselach. Z kolei zgodność z OpenCV
potwierdza, że nie jest to specyficzna, niestandardowa konwencja jednej biblioteki, ale część
szerszego, powszechnie używanego modelu kamery otworkowej z dystorsją.

\section{Wnioski}
Funkcja \texttt{intrinsics\_to\_dict(intr)} zapisuje standardowy zestaw parametrów
wewnętrznych kamery zgodny z definicją struktury \texttt{rs2\_intrinsics} w Intel RealSense
SDK 2.0. Najważniejsze wnioski interpretacyjne są następujące:
\begin{enumerate}
    \item \texttt{width} i \texttt{height} są rozmiarem obrazu w pikselach.
    \item \texttt{fx} i \texttt{fy} są ogniskowymi wyrażonymi w jednostkach pikselowych,
          a nie w milimetrach.
    \item \texttt{cx} i \texttt{cy} odpowiadają punktowi głównemu obrazu i są mierzone jako
          przesunięcie od lewego oraz górnego brzegu obrazu.
    \item \texttt{distortion.coefficients} dla modelu Brown--Conrady należy interpretować jako
          $[k_1, k_2, p_1, p_2, k_3]$.
    \item Podstawą tych stwierdzeń są oficjalna dokumentacja Intel RealSense / librealsense2
        oraz zgodna z nią dokumentacja OpenCV.
\end{enumerate}

\section*{Źródła}
\begin{enumerate}
    \item Intel RealSense / librealsense2, dokumentacja struktury \texttt{rs2\_intrinsics}: \href{https://docs.ros.org/en/rolling/p/librealsense2/generated/structrs2__intrinsics.html}{docs.ros.org/en/rolling/p/librealsense2/generated/structrs2\_\_intrinsics.html}.
    \item Intel RealSense / librealsense2, dokumentacja enumeracji \texttt{rs2\_distortion}: \href{https://docs.ros.org/en/rolling/p/librealsense2/generated/enum_rs__types_8h_1a96a8444bedb23db02c92712e0a0c37ae.html}{docs.ros.org/en/rolling/p/librealsense2/generated/enum\_rs\_\_types\_8h\_1a96a8444bedb23db02c92712e0a0c37ae.html}.
    \item OpenCV 4.x, moduł \texttt{calib3d}, dokumentacja macierzy wewnętrznej kamery i współczynników dystorsji: \href{https://docs.opencv.org/4.x/d9/d0c/group__calib3d.html}{docs.opencv.org/4.x/d9/d0c/group\_\_calib3d.html}.
\end{enumerate}

\end{document}